\section{Supported reciprocals in the actual zero-dynamics identities}
\label{sec:rd}
\newcounter{rd}\renewcommand{\therd}{RD\arabic{rd}}
\newcommand{\RD}{\refstepcounter{rd}\par\medskip\noindent\textbf{\therd.}\ }

\RD\textbf{Supported reciprocals and the exact heat source.}
This section computes the inverse operation supplied by the support
semilattice, a failure of the resulting totalized energy's monotonicity
at a collision, and the exact relation of that energy to the original
heat dynamics.

The newly supplied \path{arXiv-1801.05914v5.tar.gz} is byte-identical
to the retained Rodgers--Tao v5 original-author archive, SHA256
\path{cc3cd07ea418ec067bca6dce49ace0b4a2c67afe8d891fb82a1bdc5a1f63d6ed}.
The actual passages read for this calculation are the definitions and
proof setup, the full section \emph{Dynamics of zeroes}, the gap
propagation argument, and the full sections \emph{Controlling the energy
at time 0} and \emph{Contradicting pair correlation}. Relevant original
TeX labels are \texttt{ode}, \texttt{ident}, \texttt{tlar},
\texttt{tlar-iter}, and \texttt{picketfence}. The technical integrated
energy argument is cited as part of their theorem; it is not claimed
to have been reproduced here.

\RD\textbf{The original reciprocal, including supported zero.}
For $S=G_L(\mathbb R)$ define
\[
 \widehat a^{\dagger}=\widehat{1/a}\quad(a\ne0),
 \qquad z_\lambda^{\dagger}=z_\lambda\quad(\lambda\in L).
\]
This gives the multiplicative inverse in the inverse-monoid sense:
$xx^\dagger x=x$, $x^\dagger xx^\dagger=x^\dagger$, and
$(x^\dagger)^\dagger=x$. For nonzero amplitude these equations hold
in its ordinary multiplicative group; for a support element they are
the idempotence identities $z_\lambda^2=z_\lambda$. Uniqueness follows
as well. A nonzero-amplitude $x$ forces the amplitude and top support
of its inverse. If $x=z_\lambda$, the second inverse identity forces
the inverse to have zero amplitude, say $z_\mu$; the two identities
then give $\lambda\le\mu$ and $\mu\le\lambda$.
In particular
\[
 e^\dagger=e,\qquad ee^\dagger=e,\qquad
 \widehat a\widehat a^\dagger=\widehat1\quad(a\ne0).
\]
This acknowledges supported zero's genuine inverse in its local
multiplicative group. Replacing the middle right-hand side by
$\widehat1$ would instead perform the unital localization $S\to L$.
Both operations and their connecting maps are proved in the
semilattice chapter.

\RD\textbf{Exact transport on the distinct-root locus.}
For a finite set of labelled real roots define
$D_{jk}=x_j-x_k$ and suppose the displayed differences are nonzero.
The map $a\mapsto\widehat a$ identifies $\mathbb R$ with the top
additive group of $S$, whose identity is $e$. It preserves addition,
multiplication, the unit and additive inverses relative to $e$.
For every nonzero difference it also preserves the displayed
reciprocal. Consequently every rational identity whose denominators
are products of these differences has exactly the same lifted identity.
This follows by induction on additions, multiplications and reciprocal
operations in its expression. The projection $p:S\to\mathbb R$
is the inverse on that group. Differentiation along a top-fibre curve
$\widehat{f(t)}$ means $\widehat{f'(t)}$, so the identities commute
with differentiation too. Order on this fibre is the transported real
order, $\widehat a\le\widehat b$ exactly when $a\le b$; no order is
asserted between distinct lower support fibres.

For principal-value sums one first uses the exact finite identities,
then takes the original amplitude limit. This defines the topology and
the limit in the top fibre by the bijection $p$. It neither adds nor
removes any summand. In particular, retaining $e$ in the coefficient
semiring changes none of the source's identities on their stated
distinct-root domain.

\RD\textbf{The precise monotonicity identity, including its environment.}
Rodgers--Tao work under the contradiction assumption $\Lambda<0$;
their Theorem \texttt{dynam} states, for $\Lambda<t\le0$,
\[
 x_k'=2\sum_{j\ne k}'\frac1{x_k-x_j}.
\]
The zeros in this range are real and simple, by the cited
Csordas--Smith--Varga result. Their finite-set energy is
$E^K=\sum_{k,k'\in K,k\ne k'}(x_k-x_{k'})^{-2}$, and Lemma
\texttt{ident}(iii) states exactly
\begin{align*}
 (E^K)'={}&
 \sum_{\substack{j\notin K\\ k,k'\in K,\ k\ne k'}}
 \frac4{(x_k-x_{k'})^2(x_k-x_j)(x_{k'}-x_j)}\\
 &-2\sum_{\substack{k,k'\in K\\k\ne k'}}
 \left(\frac2{(x_k-x_{k'})^2}
 -\sum_{\substack{k''\in K\\k''\ne k,k'}}
 \frac1{(x_{k''}-x_k)(x_{k''}-x_{k'})}\right)^2.
\end{align*}
The external sum is present. A finite subsystem's energy is not
asserted to decrease without controlling it. The preceding transport
proof applies to every term, square, sign and cutoff here. In section8,
Proposition \texttt{tlar-iter} controls those external terms across
time intervals at most $1/(100\log^2T)$, with an interval shrunk by
$\log^3T$ at each end. Supported-zero inversion changes none of those
terms, because no denominator there is supported zero.

This reading also makes the logical direction exact. Their contradiction
excludes $\Lambda<0$, leaving both $\Lambda=0$ and $\Lambda>0$.
The latter is entirely compatible with their theorem and is the desired
RH-disproof direction. A failure of some new extension of energy
monotonicity, by itself, does not establish $\Lambda>0$; that inference
requires a calculation on the original arithmetic family at its time
zero, rather than on a separately extended function of colliding roots.

\RD\textbf{The full collision defect of the reciprocal substitution.}
The failure of cancellation at a collision is a computable object.
For the map from labelled roots to coefficients,
\[
 \sigma(x_1,x_2)=(a,b)=(x_1+x_2,x_1x_2),\qquad
 P(Z)=Z^2-aZ+b,
\]
backward heat gives the exact coefficient vector $(a',b')=(0,-2)$.
Away from $x_1=x_2$, the root velocities
$x_1'=2/(x_1-x_2)$, $x_2'=2/(x_2-x_1)$ give precisely that
vector after $D\sigma$.
At $(r,r)$, however,
\[
 D\sigma_{(r,r)}(v_1,v_2)=(v_1+v_2,r(v_1+v_2)).
\]
The cokernel is a line, explicitly received by
$\ell_r(\dot a,\dot b)=\dot b-r\dot a$. Its value on the actual
heat vector is $-2$. This is a nonzero obstruction to any finite
differentiable labelled-root lift at the collision.
The semilattice generalized reciprocal sets both colliding velocities
to the supported value $e$; their amplitude vector is $(0,0)$, whose
image under $D\sigma$ is $(0,0)$, leaving the same defect $(0,-2)$.
Thus the full coefficient heat vector survives, while a proposed finite
root vector fails by this explicitly calculated cokernel class.

The replacement object supplied by this obstruction is the ramified
root space $\{(h,Z):Z^2=2h\}$, where $h=Z^2/2$. Its tangent map to
time vanishes at $Z=0$, and the vector field lifting $\partial_h$ on
$Z\ne0$ is $Z^{-1}\partial_Z$. It has a simple pole in this
coordinate. Both the ramification and the pole are calculated exactly;
neither is silently removed by assigning a finite reciprocal at zero.

\RD\textbf{A real loss of monotonicity after totalizing energy.}
For the same actual polynomial heat solution $Q_h(Z)=Z^2-2h$ and
$h>0$, its two roots have gap $d(h)=2\sqrt{2h}$. Therefore its
ordered-pair energy is
\[
 E(h)=\frac2{d(h)^2}=\frac1{4h},\qquad E'(h)=-\frac1{4h^2}<0.
\]
At $h=0$ the gap is the supported zero $e$. If one evaluates the
same expression by the generalized inverse in RD2, the assigned
energy is $2e^2=e$, of amplitude zero. Consequently the totalized
function $E^\dagger(0)=0$, $E^\dagger(h)=1/(4h)$ for $h>0$ is
not nonincreasing on $[0,\infty)$: for every $h>0$ it exceeds its
value at zero. This proves a precise version of the user's proposed
monotonicity failure.

It also determines its logical consequence. Every zero of $Q_h$ is
real for every $h\ge0$, including the double zero at zero. Hence this
failure of the extended energy's monotonicity is compatible with
real-rootedness at time zero. It cannot serve as the implication from
failure of that extended energy law to failure of the real-zero
condition. The original energy has $\lim_{h\downarrow0}E(h)=+\infty$;
the generalized-inverse value is a different, discontinuous extension,
and the two are related by the explicit reciprocal and amplitude maps.

The same distinction is visible without orders or inequalities. In
the ramification coordinate $q:=Z$ one has $h=q^2/2$ and
$E=1/(2q^2)$. Its Laurent germ has valuation $-2$ and leading
coefficient $1/2$. The valued germ field $\mathbb C((q))$, with its
valuation and residue of $q^2E$, retains both data; evaluation of
regular functions at $q=0$ does not extend to it as a nonzero
ring homomorphism. The split special-fibre receiver of the localized
chart is $L$ as proved in \eqref{eq:eh-boundary}. In the original
coordinate $r=iZ/2$ the same identities are $h=-2r^2$ and
$E=-1/(8r^2)$, as computed in \eqref{eq:eh-energy-coordinates}.
That receiver retains support;
the Laurent valuation retains the pole order as additional data.
This identifies exactly which boundary information is needed if an
energy argument is to cross the collision.

\RD\textbf{The differential equation retained at supported time zero.}
At a simple zero of the actual entire function, differentiating
$H_t(x(t))=0$ gives
$x'=H_{ZZ}/H_Z$, by $\partial_tH=-H_{ZZ}$.
At a double zero $(t_0,x_0)$ where $H_{ZZ}\ne0$, this equation in
undivided form is
\[
 H_Z x'=H_{ZZ}.
\]
Its value there has left amplitude zero for every finite velocity,
and nonzero right amplitude. In the top split fibre it is the
unequal pair $e$ and $\widehat{H_{ZZ}}$. Making $e$ a unit by
localization sends both to the top support and thus solves only
the support image of the equation. Applying the amplitude map
before localization retains the displayed nonzero differential
defect. These are exact maps between the two equations; no
arithmetic sign or nonreal zero is obtained by their identification.

\RD\textbf{An actual arithmetic sign at the origin.}
The moments in \eqref{eq:eh-moments} evaluate a sign without a finite-family
replacement. For every real $t$ put
$M_j(t)=\int_0^\infty u^{2j}e^{tu^2}\Phi(u)\,du>0$.
The complete original source gives
\[
 H_t(0)=M_0(t)>0,\quad H_t'(0)=0,\quad H_t''(0)=-M_1(t)<0,
\]
and hence
\[
 (H_t'(0))^2-H_t(0)H_t''(0)=M_0(t)M_1(t)>0.
\]
All derivatives and signs follow from the domination and positivity
proved with \eqref{eq:eh-moments}. Thus neither spatial zero $Z=0$ nor this curvature
test at that spatial point supplies a negative witness at time $t=0$.
The spatial point and the time parameter are retained separately.
No inference about the sign at other spatial points is made.

\RD\textbf{The heat shift, its time-zero fibre, and the full support lattice.}
This calculation tests the inverse-monoid operation while retaining
the distinguished time-zero fibre. Write $c_{\rm h}\in\mathbb R$
for the prescribed heat shift $c$, to distinguish it from the third
target coordinate $c$ of the original polynomial $F$. For every
value of this retained parameter define
\begin{equation}
 Q_t^{[c_{\rm h}]}(Z)=Z^2-2(t-c_{\rm h})
 =Z^2-2t+2c_{\rm h},\qquad
 \partial_tQ_t^{[c_{\rm h}]}=-\partial_Z^2Q_t^{[c_{\rm h}]}=-2.
 \label{eq:rd-shifted-heat}
\end{equation}
In this test family $t$ marks its time-zero fibre. Its use as a
coordinate does not assert equality with the time of the particular
arithmetic function $H_t$.

First keep $c_{\rm h}$ indeterminate. There is an exact isomorphism
of $\mathbb R[c_{\rm h}]$-algebras
\begin{equation}
 \theta:\mathbb R[c_{\rm h},h,Z]\longrightarrow
              \mathbb R[c_{\rm h},t,Z],\qquad
 \theta(h)=t-c_{\rm h},\quad\theta(Z)=Z,
 \qquad \theta^{-1}(t)=h+c_{\rm h}.
 \label{eq:rd-shifted-coordinate-map}
\end{equation}
The generator formulas compose to the identity and hence prove the
isomorphism. They send $Z^2-2h$ to
$Q_t^{[c_{\rm h}]}$ and induce an isomorphism of the quotient root
algebras. The two distinguished evaluations satisfy
\begin{equation}
 \operatorname{ev}_{t=0}\circ\theta
   =\operatorname{ev}_{h=-c_{\rm h}},\qquad
 \operatorname{ev}_{t=c_{\rm h}}\circ\theta
   =\operatorname{ev}_{h=0}.
 \label{eq:rd-shifted-marked-fibres}
\end{equation}
Thus the isomorphism retains the marked fibre by specifying its
image; it does not send time zero to the collision unless
$c_{\rm h}=0$. Parameter specialization to any real $c_{\rm h}$
commutes with these maps.

Here are the heat maps in their exact receiving category. Let
$A=\mathbb R[c_{\rm h},t,Z]$. For a real increment $\delta$ set
\begin{equation}
 \mathscr H_\delta f
   =\sum_{j\ge0}\frac{(-\delta)^j}{j!}\partial_Z^{2j}f,
 \qquad
 \mathscr H_\delta Q_t^{[c_{\rm h}]}
   =Q_{t+\delta}^{[c_{\rm h}]}.
 \label{eq:rd-shifted-heat-operator}
\end{equation}
The sum is finite on each polynomial. This is an additive
$\mathbb R[c_{\rm h},t]$-linear map. Multiplying the two finite
operator sums and collecting $\partial_Z^{2n}$ gives the coefficient
$\sum_{j=0}^n(-\delta)^j(-\epsilon)^{n-j}/(j!(n-j)!)
=(-\delta-\epsilon)^n/n!$. Consequently
$\mathscr H_\delta\mathscr H_\epsilon
=\mathscr H_{\delta+\epsilon}$, and
$\mathscr H_{-\delta}$ is its inverse under composition.
The second equality in \eqref{eq:rd-shifted-heat-operator} follows
because $\partial_Z^2Q=2$ and all higher even derivatives vanish.
In particular $Q_t^{[c_{\rm h}]}$ is exactly the formal heat image
$\exp(-t\partial_Z^2)(Z^2+2c_{\rm h})$.

For every original bounded distributive lattice $L$ retain the carrier
\[
 G_L(A)=\{(0,\lambda):\lambda\in L\}
          \cup\{(f,1_L):f\in A\},\quad
 e=(0,1_L),\quad\tau=(0,0_L).
\]
Its addition and multiplication are respectively
$(f,\lambda)+(g,\mu)=(f+g,\lambda\vee\mu)$ and
$(f,\lambda)(g,\mu)=(fg,\lambda\wedge\mu)$.
The lifted heat map is
\begin{equation}
 \widetilde{\mathscr H}_\delta(f,\lambda)
       =(\mathscr H_\delta f,\lambda),\qquad
 p\widetilde{\mathscr H}_\delta=\mathscr H_\delta p,
 \quad\chi_L\widetilde{\mathscr H}_\delta=\chi_L.
 \label{eq:rd-shifted-split-heat}
\end{equation}
It is well defined because any input with support below $1_L$ has
amplitude zero, which the additive map sends to zero. The displayed
addition formula proves additivity, and scalar linearity proves
$G_L(\mathbb R[c_{\rm h},t])$-linearity. It fixes every $z_\lambda$,
including $e$, and its inverse is
$\widetilde{\mathscr H}_{-\delta}$. These are automorphisms of
additive semimodules. They are not asserted to be semiring maps:
for $\delta\ne0$, $\mathscr H_\delta(Z)^2=Z^2$ whereas
$\mathscr H_\delta(Z^2)=Z^2-2\delta$.

In contrast, the coordinate isomorphism $\theta$, parameter
specialization, and both evaluations in
\eqref{eq:rd-shifted-marked-fibres} are ring homomorphisms. Each
lifts to the unital semiring homomorphism
$(f,\lambda)\mapsto(\varphi(f),\lambda)$, retaining every support
label. This formula proves all corresponding squares commute.
The joint map $(p,\chi_L)$ is injective, so these lifted maps retain
exactly the amplitude data of the stated ring maps together with the
whole original support lattice.

Now fix a real $c_{\rm h}$. The root algebra and its time-zero fibre
are
\begin{equation}
 B_{c_{\rm h}}=\mathbb R[t,Z]/(Z^2-2(t-c_{\rm h})),\qquad
 B_{c_{\rm h}}/(t)
   =\mathbb R[Z]/(Z^2+2c_{\rm h}).
 \label{eq:rd-shifted-root-fibre}
\end{equation}
The lifted quotient $G_L(B_{c_{\rm h}})\to
G_L(B_{c_{\rm h}}/(t))$ identifies precisely pairs with the same
support label and amplitude difference in $tB_{c_{\rm h}}$.
Indeed this follows separately from the two coordinates of its
formula, and proves the asserted kernel congruence in both
directions. It does not erase a support label or replace a supported
zero by $\tau$.

The real algebras in \eqref{eq:rd-shifted-root-fibre} can be
identified completely. If $c_{\rm h}<0$, write
$a=\sqrt{-2c_{\rm h}}>0$. Evaluation at $a$ and $-a$ gives an
isomorphism with $\mathbb R\times\mathbb R$: for any desired
values $u,v$, the degree-one polynomial
$(u+v)/2+(u-v)Z/(2a)$ has those values, and the kernel is
$(Z^2-a^2)$ by division. If $c_{\rm h}=0$ the algebra is
$\mathbb R[\eta]/(\eta^2)$, with $Z\mapsto\eta$. If
$c_{\rm h}>0$, evaluation $Z\mapsto i\sqrt{2c_{\rm h}}$ gives
an isomorphism with $\mathbb C$ as a real algebra: the images of
$1,Z$ form a real basis, and the relation is satisfied. These
isomorphisms lift under $G_L$ as before. They calculate the entire
time-zero fibre, including its nilpotent when it is singular.

At the collision time $t=c_{\rm h}$ the fibre is always
$\mathbb R[\eta]/(\eta^2)$. This agrees with the exact
ramification coordinate $q=Z$, $t=c_{\rm h}+q^2/2$. The
point-value receiver $G_L(\mathbb C)$ is the same inverse monoid
for every $c_{\rm h}$, with the inverse of RD2 and $e^\dagger=e$.
The collision algebra also retains its nonzero nilpotent $\eta$;
that element does not acquire a generalized inverse in the
algebra. Indeed $\eta b\eta=0\ne\eta$ for every $b$, so the
first inverse identity fails. After applying $G_L$ the same
amplitude calculation proves this failure for $\widehat\eta$.
Thus the point-value inverse operation and the full singular fibre
are related by their actual evaluation map $\eta\mapsto0$;
the multiplicity has not been discarded.

For real $t$ the roots themselves are
\begin{equation}
 \begin{cases}
  \displaystyle\pm i\sqrt{2(c_{\rm h}-t)},&t<c_{\rm h},\\
  0\text{ with multiplicity two},&t=c_{\rm h},\\
  \displaystyle\pm\sqrt{2(t-c_{\rm h})},&t>c_{\rm h}.
 \end{cases}
 \qquad
 \{t:Q_t^{[c_{\rm h}]}\text{ has only real roots}\}
       =[c_{\rm h},\infty).
 \label{eq:rd-shifted-threshold}
\end{equation}
This proves that the threshold is exactly $c_{\rm h}$ and proves
its behaviour at equality. At the marked time zero the roots are
nonreal for $c_{\rm h}>0$, a real double root for $c_{\rm h}=0$,
and distinct real roots for $c_{\rm h}<0$. The evaluation map
\[
 G_L(\mathbb C[c_{\rm h},t,Z])\longrightarrow G_L(\mathbb C),
 \qquad(f,\lambda)\longmapsto
       (f(c_{\rm h},t,Z),\lambda)
\]
detects each root by the equality
$\widehat{Q_t^{[c_{\rm h}]}(Z)}=e$. Composing with $\chi_L$
sends that value, and every supported nonzero value, to $1_L$.
This gives the exact map that forgets the root condition; the joint
receiver $(p,\chi_L)$ retains it.

Finally, on $t>c_{\rm h}$ the pair gap is
$d=2\sqrt{2(t-c_{\rm h})}$ and its ordered-pair energy is
\begin{equation}
 E_{c_{\rm h}}(t)=\frac1{4(t-c_{\rm h})},\qquad
 E_{c_{\rm h}}'(t)=-\frac1{4(t-c_{\rm h})^2},\qquad
 E_{c_{\rm h}}^\dagger(c_{\rm h})=0.
 \label{eq:rd-shifted-energy}
\end{equation}
The last term is the amplitude of the supported value
$\widehat2(d^\dagger)^2=e$ at $d=e$; it is not the limit of the
first term. Therefore the same failure of nonincreasing behaviour
on $[c_{\rm h},\infty)$ occurs for every shift. The root-to-coefficient
map still has heat vector $(0,-2)$, and at the collision $(0,0)$
its differential has image $\{(v,0):v\in\mathbb R\}$.
The cokernel functional is $\ell_0(\dot a,\dot b)=\dot b$, so
the unresolved finite-velocity class is still exactly $-2$.
Time translation changes none of these local identities.

The same full lattice $L$, the same point-value inverse monoid,
and the same genuine local inverse $e^\dagger=e$ therefore occur
with every real threshold $c_{\rm h}$. In particular, both
$c_{\rm h}=0$ and $c_{\rm h}>0$ occur, so these inverse and
support data alone do not choose between zero and a positive
threshold, even after restricting attention to nonnegative
thresholds. The singularity at the marked time zero is retained
explicitly by \eqref{eq:rd-shifted-root-fibre}; it is present
precisely when $c_{\rm h}=0$. This is a test of the proposed
inference from the inverse operation. The initial amplitudes
$Z^2+2c_{\rm h}$ are not the arithmetic amplitudes $H_0$, whose
moments and singularity question remain unchanged. No conclusion
about that arithmetic singularity or RH is obtained by replacing
those amplitudes with this family.
